\subsection{A variable-group quotient of the actual path correspondence}

The actual quotient of a space by its group action
is a different operation from the coefficient-object
inversion excluded in \fid{F68}.
Here a global quotient normalizes the whole
\fid{F62} three-point path correspondence, including
its diagonal degeneration. The four-point test
then isolates a precise additional coherence problem.

Let \(k\) be any field and put
\[
 Y=\mathbb P^1_x\times\mathbb P^1_z,\qquad
 Z=\mathbb P^1_t\times Y,\qquad p:Z\to Y.
\]
Let \(\Gamma_x,\Gamma_z\) be the section divisors
\(t=x,t=z\). Define their Cartier sum and its algebra by
\begin{equation}\label{eq:variableQuotientDivisor}
 \mathscr D=\Gamma_x+\Gamma_z,\qquad
 I_{\mathscr D}=I_{\Gamma_x}I_{\Gamma_z},\qquad
 \pi=p|_{\mathscr D},\qquad
 \mathcal A=\pi_*\mathcal O_{\mathscr D}.
\end{equation}
The Cartier sum retains a double point when \(x=z\).
Write \(\Delta\subset Y\) for that diagonal.

\begin{finding}[\fid{F69}, actual variable-group normalization]
\label{f:variableGroupQuotient}
The morphism \(\pi\) is finite locally free of degree two.
The group of its units modulo scalar units is represented
by a smooth affine group scheme
\[
 \mathcal J=(\operatorname{Res}_{\mathscr D/Y}\mathbb G_m)/
                         \mathbb G_m
\]
over \(Y\). Its geometric fibers are \(\mathbb G_m\)
off \(\Delta\) and \(\mathbb G_a\) on \(\Delta\).

The actual path scheme
\[
 \mathcal P=\{(x,y,z):x\ne y,\ y\ne z\}
\]
has an explicit \(\mathcal J\)-action making
\(\mathcal P\to Y\) a Zariski locally trivial torsor.
Consequently
\[
                         [\mathcal P/\mathcal J]\simeq Y
\]
as fppf stacks, with no residual inertia.
The construction retains the nonzero off-diagonal
Chern extension of \fid{F62}.

On the four-point path space
\(\mathcal W=\{x\ne y,\ y\ne z,\ z\ne w\}\),
the second named group family \(\mathcal J_{y,w}\)
cannot descend equivariantly under the first action
\(\mathcal J_{x,z}\).
This is a failure of those specified transport data,
not an exclusion of all weak or higher coherent
stack normalizations.
\end{finding}

\begin{proof}
\par\smallskip\noindent\textbf{The finite-flat algebra.}\quad
For \(a\in\{0,1,\infty\}\), put
\[
 U_a=(\mathbb P^1\setminus\{a\})_x
                \times(\mathbb P^1\setminus\{a\})_z.
\]
These opens cover \(Y\): two geometric endpoints
cannot contain all three specified points.
On \(U_a\), the divisor avoids \(t=a\).
In the associated affine coordinate it has
the monic equation \((t-x)(t-z)=0\).
Writing \(h=z-x\) and \(c=t-x\) gives
\begin{equation}\label{eq:variableQuotientAlgebra}
 \mathcal A|_{U_a}=R[c]/(c^2-hc),\qquad
 \operatorname{ev}_x(c)=0,\quad
 \operatorname{ev}_z(c)=h.
\end{equation}
It is free on \(1,c\), proving finiteness and
local freeness of degree two directly.

The section ideals are invertible, and
\[
 I_{\Gamma_x}/I_{\Gamma_x}I_{\Gamma_z}
       \simeq I_{\Gamma_x}|_{\Gamma_z}
       \simeq\mathcal O_Y(-\Delta).
\]
Thus, with \(\mathcal L_0=\mathcal O_Y(-\Delta)\),
evaluation at \(x\) and the unit give a split sequence
\begin{equation}\label{eq:variableQuotientAugmentation}
 0\longrightarrow\mathcal L_0
 \longrightarrow\mathcal A
 \xrightarrow{\operatorname{ev}_x}\mathcal O_Y
 \longrightarrow0.
\end{equation}
The kernel multiplication is determined by the
natural inclusion \(\mathcal L_0\to\mathcal O_Y\);
locally it is \(c^2=hc\).
Off the diagonal the geometric algebra is
\(\bar k\times\bar k\); on it the algebra is
\(\bar k[\epsilon]/(\epsilon^2)\).
The nilpotent is not discarded.
The monic presentation also verifies arbitrary
base-change compatibility.

\par\smallskip\noindent\textbf{The actual quotient group.}\quad
Let \(\mathcal G=\operatorname{Res}_{\mathscr D/Y}\mathbb G_m\).
It is explicitly the open of units in the rank-two
vector bundle \(\mathcal A\). In the local algebra,
\[
                 N(a+bc)=a(a+hb)
\]
is the determinant of multiplication, so this
unit open is represented by its nonvanishing.
Evaluation at \(x\) is a split group homomorphism
\(\mathcal G\to\mathbb G_m\).
Every unit has an invertible value there; division
by that scalar gives a canonical product decomposition
\[
 \mathcal G\simeq\mathbb G_m\times\mathcal J,\qquad
 \mathcal J=\ker(\operatorname{ev}_x).
\]
In particular the scalar fppf quotient is represented.

Globally \(\mathcal J\) is an open in the total space
of \(\mathcal L_0\). Locally its elements are \(1+uc\):
\begin{equation}\label{eq:variableQuotientGroupLaw}
 \mathcal J|_{U_a}
   =\operatorname{Spec}R[u,(1+hu)^{-1}],
 \qquad u\odot v=u+v+huv.
\end{equation}
The identity is zero and the inverse is
\(-u/(1+hu)\). Associativity and commutativity
are actual multiplication in \(\mathcal A\);
the formulas glue by the kernel construction.
This is a smooth affine group of relative dimension one.

When \(h\) is invertible, the coordinate \(1+hu\)
identifies the fiber with \(\mathbb G_m\);
when \(h=0\), the group law is additive.
On the whole diagonal the group is the additive
vector group of
\(\mathcal L_0|_\Delta=\mathcal O_{\mathbb P^1}(-2)\).
No globally constant additive bundle is asserted.
Off \(\Delta\), the two evaluations give a global
multiplicative coordinate, for example
\(\operatorname{ev}_x/\operatorname{ev}_z\).

\par\smallskip\noindent\textbf{Evaluation and the path scheme.}\quad
Let \(H=H^0(\mathbb P^1,\mathcal O(1))\), and put
\(\mathcal M=\pi_*(\mathcal O(1)|_{\mathscr D})\).
Evaluation gives
\begin{equation}\label{eq:variableQuotientEvaluation}
                 H\otimes_k\mathcal O_Y
                       \xrightarrow{\sim}\mathcal M.
\end{equation}
Indeed, on the preceding affine charts, the sections
\(1,t\) restrict to a basis \(1,c\), with \(t=x+c\)
and determinant one.
Equivalently,
\[
 0\to\mathcal O_Z(-1)\otimes p^*\mathcal O_Y(-1,-1)
   \to\mathcal O_Z(1)\to\mathcal O_{\mathscr D}(1)\to0
\]
and the cohomology vanishing of \(\mathcal O_{\mathbb P^1}(-1)\)
give the same result.
The direct local determinant proof includes the double point.

Let \(\mathcal T_{\rm fr}\subset\mathcal M\) be the
open of \(\mathcal A\)-generators:
\[
 \mathcal T_{\rm fr}
  =\operatorname{Res}_{\mathscr D/Y}
       \operatorname{Isom}_{\mathscr D}
       (\mathcal O_{\mathscr D},\mathcal O(1)|_{\mathscr D}).
\]
The condition that \(a\mapsto as\) is an isomorphism
from \(\mathcal A\) to \(\mathcal M\) is a determinant
nonvanishing condition.
Projectivizing these frames gives an open scheme
\[
 \mathcal T_{\rm fr}/\mathbb G_m
       \subset\mathbb P_{\rm lines}(\mathcal M).
\]
Here projective space parametrizes lines of sections.
The scalar quotient is the usual projective-bundle
quotient, verified on local scalar frame charts.

A nonzero section of \(\mathcal O(1)\), up to scalar,
has a unique zero \(y\).
It generates \(\mathcal O(1)|_{\mathscr D}\) precisely
when \(y\ne x,z\).
On a double point this is the unit criterion in
the local dual-number algebra.
The same determinant-open condition holds in families.
Consequently the zero-of-section map gives the
actual scheme identification
\begin{equation}\label{eq:variableQuotientActualPath}
 \mathcal T_{\rm fr}/\mathbb G_m
       \simeq\mathcal P
       =\{(x,y,z):x\ne y,\ y\ne z\}.
\end{equation}

\par\smallskip\noindent\textbf{The action and torsor maps.}\quad
Units of \(\mathcal A\) act by multiplication on
\(\mathcal M\). Scalar units act by scalar multiplication,
so this induces an actual \(\mathcal J\)-action on
its projectivization and on \(\mathcal P\).
For two frames \(s,t\), their ratio \(t/s\) is a
unique unit. Replacing either by a scalar multiple
changes the ratio by a scalar unit.
Thus the ratio defines the inverse of
\begin{equation}\label{eq:variableQuotientTorsorMap}
 \mathcal J\times_Y\mathcal P
    \longrightarrow\mathcal P\times_Y\mathcal P,
             \qquad(g,p)\longmapsto(gp,p).
\end{equation}
Local frame ratios are regular and scalar-independent,
so they glue to an actual inverse.

There are Zariski local sections: on \(U_a\) use
the section of \(\mathcal O(1)\) vanishing at \(a\),
or equivalently the path point \(y=a\).
It is a generator since neither endpoint equals \(a\).
Hence \(\mathcal P\to Y\) is a Zariski locally
trivial \(\mathcal J\)-torsor.

In a local \(\mathcal A\)-frame of \(\mathcal M\),
multiplication by \(1+uc\) acts by
\begin{equation}\label{eq:variableQuotientProjectiveAction}
                  [a:b]\longmapsto[a:b+ua+hub].
\end{equation}
The determinant is \(1+hu\), which is invertible.
This gives an action on the entire projective bundle.
It fixes the two endpoint zero-lines, with the
corresponding coincident fixed line on the diagonal.
Its generator-open restriction is the torsor action.

\par\smallskip\noindent\textbf{The quotient stack.}\quad
An object of \([\mathcal P/\mathcal J]\) over a
scheme \(T\to Y\) is a \(\mathcal J_T\)-torsor \(Q\)
and an equivariant map \(Q\to\mathcal P_T\).
An equivariant map between torsors is an isomorphism:
locally both are the group and the map is translation.
The object is therefore uniquely isomorphic to
\((\mathcal P_T,\operatorname{id})\), with no
nonidentity automorphism preserving that map.
This proves canonically
\begin{equation}\label{eq:variableQuotientIsBase}
                          [\mathcal P/\mathcal J]\simeq Y.
\end{equation}
The equivalence commutes with base change.
The quotient-stack/torsor description is the known
Stacks Proposition~95.15.3, tag~04WM
\cite{StacksVariableQuotients}; all its torsor
hypotheses have been proved explicitly here.

\par\smallskip\noindent\textbf{The exact finite-field fibers.}\quad
Over \(\mathbb F_Q\), distinct endpoints give
\(\mathcal A=\mathbb F_Q^2\), a group
\(\mathbb G_m\) of order \(Q-1\), and \(Q-1\)
path points. Equal endpoints give the dual numbers,
a group \(\mathbb G_a\) of order \(Q\), and \(Q\)
path points. In both cases the quotient fiber is
one point with trivial automorphism.
The whole path count is
\[
 Q(Q+1)(Q-1)+(Q+1)Q=Q^2(Q+1);
\]
the quotient has \((Q+1)^2\) points.
These counts follow the actual free transitive action;
they were not used to infer that action.

\par\smallskip\noindent\textbf{The retained Chern extension.}\quad
On \(O=Y\setminus\Delta\), the tautological lines
\(V_x,V_z\) split \(k^2\).
A line \(V_y\) avoiding both is the graph of a
nonzero map \(V_x\to V_z\). Hence
\begin{equation}\label{eq:variableQuotientChernTorsor}
 \mathcal P|_O
       \simeq\operatorname{Hom}(V_x,V_z)^\times
       \simeq\mathcal O(1,-1)|_O^\times.
\end{equation}
Use the coordinate \(\operatorname{ev}_x/\operatorname{ev}_z\)
on \(\mathcal J|_O\).
Multiplication of a section by a unit of endpoint
values \(u_x,u_z\) multiplies its graph parameter by
\(u_x/u_z\), so the line-bundle orientation matches
the actual \fid{F62} torsor.

For this calculation take \(\ell\) invertible in \(k\),
\(E/\mathbb Q_\ell\) finite, and geometric cohomology
after base change to an algebraic closure.
Over finite fields retain the Weil action as in \fid{F62}.
The diagonal Gysin sequence gives
\[
 H^2(O_{\bar k},E(1))
   =(Eh_x\oplus Eh_z)/E(h_x+h_z).
\]
Thus \(c_1(\mathcal O(1,-1))=h_x-h_z\) is
nonzero, becoming \(2h_x\) in the quotient.
The already proved \fid{F62} triangle remains
\[
 Rf_!E\longrightarrow E_O(-1)[-2]
  \xrightarrow{c_1(\mathcal O(1,-1))}E_O
  \longrightarrow Rf_!E[1].
\]
Its nonzero geometric Chern class proves nonsplitting
after geometric base change, hence also before it.
The quotient uses the actual action; it neither
splits this extension nor tensor-inverts the
ordinary compact-cohomology object.

\par\smallskip\noindent\textbf{Four-point non-descent for the named groups.}\quad
On
\[
 \mathcal W=\{(x,y,z,w):x\ne y,\ y\ne z,\ z\ne w\},
\]
the first group \(\mathcal J_{x,z}\) acts on \(y\)
and fixes \(x,z,w\).
The second named family is \(\mathcal J_{y,w}\).
Suppose the latter had an equivariant group-family
descent datum under the first action.

Pass to an algebraic closure and fix
\(x=0,z=\infty,w=1\).
The first group is \(\mathbb G_m\), transitive on
the permitted \(y\in\mathbb G_m\).
It can send \(y=1\) to \(y=b\) for
\(b\notin\{0,1\}\); such \(b\) exists also in
characteristic two after algebraic closure.
The second group is \(\mathbb G_a\) at the first
path and \(\mathbb G_m\) at its translate.
An equivariant transport would identify these fibers
as group schemes.
They are not even isomorphic as varieties:
the coordinate ring of \(\mathbb A^1\) has only
constant units, while that of \(\mathbb G_m\)
has a nonconstant unit.
This contradiction proves the claimed non-descent,
before any cocycle condition is considered.
The initial path \((0,1,\infty,1)\) is a valid path;
its nonadjacent coincidence cannot simply be deleted
while preserving this action.
\end{proof}

\subsubsection{What a new second-stage endpoint action retains}

The first quotient nevertheless exists. Pulling
back the proved torsor identity gives
\[
 [\mathcal W/\mathcal J_{x,z}]
       \simeq R_z=\{(x,z,w):z\ne w\}.
\]
It introduces no condition \(z\ne x\).
One may define a \emph{new} second-stage action
using the endpoint group \(\mathcal J_{x,w}\).
The projective action
\eqref{eq:variableQuotientProjectiveAction}
preserves this open because it fixes \(w\).
This is not a descent of the original second family.

For \(x\ne w\), identify
\(\mathbb P^1\setminus\{w\}\) with \(\mathbb A^1\)
and \(x\) with zero.
The group is \(\mathbb G_m\), acting by a multiplicative
coordinate or its inverse.
The quotient is therefore
\[
                         [\mathbb A^1/\mathbb G_m].
\]
Its nonzero free orbit has quotient a point;
the fixed point \(z=x\) contributes \(B\mathbb G_m\).
For \(x=w\), the group is instead \(\mathbb G_a\)
acting by translations on \(\mathbb P^1\setminus\{x\}\);
that quotient is a point.

In the off-diagonal finite-field fiber the mass is
\[
                  1+\frac1{Q-1}=\frac Q{Q-1},
\]
not one. The fixed-point term has this mass because
a line over a field is trivial and its automorphism
group has \(Q-1\) elements.
Thus this particular successive quotient does not
return the endpoint scheme \(Y\).
The reverse parenthesization has the analogous
one-punctured-line issue.
Neither calculation excludes all possible
stack-valued associators or higher normalizations.

\paragraph{The finite-Weil boundary for this exact residual mass.}
For distinct endpoints over \(\mathbb F_Q\),
the second-stage fiber mass at extension degree \(n\) is
\begin{equation}\label{eq:variableQuotientResidualMass}
       s_n=\frac{Q^n}{Q^n-1}
                   =\sum_{j\ge0}(Q^{-j})^n.
\end{equation}
It cannot agree eventually with a finite
uninserted Weil Euler trace.
To prove this directly, for \(N\ge1\) and any \(r\)
form the rational Hankel matrix
\[
             H_{ab}=s_{N+a+b},\qquad 0\le a,b\le r.
\]
For a nonzero real vector \(v\), put
\(p(t)=\sum_av_at^a\). Then
\[
 v^{\mathsf t}Hv
     =\sum_{j\ge0}Q^{-jN}p(Q^{-j})^2>0.
\]
The series converges, and a nonzero polynomial
cannot vanish at all the infinitely many distinct
\(Q^{-j}\).
Thus every such determinant is a nonzero rational number.
A sequence \(\sum_{i=1}^rc_i\alpha_i^n\), however,
has this matrix of rank at most \(r\) over its
characteristic-zero coefficient field.
Taking \(N\) beyond an eventual-equality threshold
contradicts the same nonzero rational determinant.
A finite uninserted Weil Euler trace has exactly
this finite-sum form, including signed multiplicities
and nonsemisimple Frobenius.
This is the finite-trace boundary of \fid{F26}
for the particular mass above, not an exclusion
of the actual quotient stack or every higher normalization.

\paragraph{The exact additional geometry.}
The three-point quotient is an actual global
normalization, with its nilpotent collision, action,
base change and finite-field fibers proved.
Already \([\mathbb G_m/\mathbb G_m]\simeq\operatorname{Spec}k\)
uses an equivariant space; it does not invert
\(R\Gamma_c(\mathbb G_m,E)\) by an additive strong
monoidal coefficient functor.
It is therefore consistent with \fid{F68}.

The construction supplies no named nonstandard
generator relations, upper/lower lattices,
canonical labels, pure-coefficient comparison
or pointwise sign bound.
Using it coherently in a stronger correspondence
category remains additional work.
The proved four-point obstruction is exactly the
failure of the specified equivariant transport,
and the explicit alternative has the retained
fixed-point stratum just computed.

% Separate paragraph for Part III.
\paragraph{The \fid{F69} quotient in Grothendieck's language.}
The divisor ideals, rank-two algebra, invertible
\(\mathcal A\)-module and evaluation map are actual
affine and coherent-sheaf data.
EGA~I, Theorem~1.3.7, p.~86, Corollary~1.3.8,
p.~87 and Theorem~1.4.1, pp.~90--92 give the
verified module/sheaf dictionary \cite{EGA}.
Stacks Definition~29.49.1, tag~02KA fixes
the finite-locally-free convention.
Section~97.11, tag~05Y8 and Lemma~97.11.2,
tag~05YC give restriction of scalars and its
base change; Situation~95.15.1 and
Proposition~95.15.3, tag~04WM give the torsor
description of the quotient
\cite{StacksVariableQuotients}.
Representability here was proved by the norm open,
not merely inferred from a general quotient theorem.

Off the diagonal, the split multiplicative group
also fits SGA~3, I, 4.7.2--3 and cote~90,
scan~17 \cite{SGA3,Archive}.
That diagonalizable-group language is not applied
to the additive collision fibers.
The retained Chern calculation uses the exact
\fid{F62} Gysin/projective inputs:
\cite[Theorem~16.1 and Corollary~16.2, p.~108;
Theorem~23.2, pp.~139--140]{Milne}.
SGA~6, V, \S\S2--3 and cote~133, scan~18
remain the actual/virtual coefficient framework
\cite{SGA6,Archive}, with the recorded source-access
qualification; they do not supply this quotient
or a noncommutative \(\lambda\)-ring.
The new scheme construction and four-point
deductions have their full proofs above.
